\chapter{Methodik}\label{ch:methodik}

In diesem Kapitel wird das methodische Vorgehen zur Entwicklung einer LLM-gestützten Pipeline zur Spezifikationsgenerierung für das Vivian Framework beschrieben. Ziel ist es ein System zu entwerfen, welches semantisch sowie strukturell korrekte Spezifikationen automatisiert erzeugt, Benutzer*innen-Feedback entgegennimmt sowie aufkommende Fehler iterativ und eigenständig behebt.

Diese Arbeit ist in Kooperation mit dem Institut für angewandte Informatik (IFAI) der Technischen Hochschule Nürnberg Georg-Simon-Ohm entstanden. Aufgrund dieser Zusammenarbeit sowie der Expertise im \ref{sec:vivian_framework} beschriebenen Vivian Framework, wurde dieses als Grundlage zur Umsetzung einer solchen Pipeline gewählt. Da dieses Framework interaktive Prototypen ausschließlich über eine maschinenlesbare Beschreibung des Verhaltens möglich macht, bietet es eine Grundlage, auf der LLM-gestützte Verfahren aufsetzen können.
Im Folgenden wird die LLM-gestützte Pipeline zur Generierung von Vivian Spezifikationen \textit{Vivian Specification Generator} (VSG) genannt.



\section{Forschungs- und Entwicklungsansatz}



\section{Anforderungsanalyse}


Im Mittelpunkt steht die Analyse vorhandener Szenendaten, die Wiederverwendung vorhandener Geometrien, die Erzeugung Vivian-kompatibler Spezifikationen, die niedrigschwellige Eingabe gewünschter Interaktionen sowie die frühzeitige Erkennung und Korrektur syntaktischer, semantischer sowie referenzieller Fehler.
Zur Nachvollziehbarkeit wurden die Anforderungen in funktionale Anforderungen (A), nicht-funktionale Anforderungen (N) sowie domänenspezifische Anforderungen (D) unterteilt.


\subsection*{Funktionale Anforderungen}

\textbf{A1:} Der VSG muss aus bestehenden Szenenrepräsentationen und natürlichsprachlichen Interaktionsbeschreibungen Vivian-Spezifikationsdateien erzeugen, die alle für die beschriebene Interaktion notwendigen Interaktionselemente, Visualisierungselemente, Zustände und Übergänge enthalten.

\textbf{A2:} Relevante Objekte und Teilobjekte müssen vom VSG explizit identifiziert und Rollen zugeordnet werden, sodass diese eindeutig über Namensreferenzen in den Vivian-Spezifikationen referenziert werden können.

\textbf{A3:} Der VSG muss erzeugte Spezifikationen syntaktisch, semantisch und referenziell prüfen.

\textbf{A4:} Der VSG muss im Generierungsprozess erkannte Fehler in einem begrenzten Korrekturprozess erneut an die zuständigen Verarbeitungsschritte zurückführen.

\textbf{A5:} Um Szeneninterpretationen überprüfen und gegebenenfalls korrigieren zu können, muss ein Human-in-the-Loop-Schritt implementiert werden, um Nutzer*innen die Möglichkeit zur Prüfung und Korrektur des Interaktionsplans zu geben, bevor aus diesem eine Spezifikation erzeugt wird.

\textbf{A6:} Der VSG muss Nutzer*innen ermöglichen, gewünschte Interaktionen natürlichsprachlich zu beschreiben, ohne dass Vivian-Spezifikationen manuell erstellt oder bearbeitet werden müssen.

\subsection*{Nicht-funktionale Anforderungen}

\textbf{N1:} Der VSG muss als modulare Pipeline mit definierten Ein- und Ausgabeformaten pro Modul aufgebaut sein, sodass einzelne Verarbeitungsschritte unabhängig beschrieben, validiert und erweitert werden können.

\textbf{N2:} Zwischenartefakte und Entscheidungen des VSG müssen protokolliert und abgespeichert werden, um die Nachvollziehbarkeit zu verbessern. Dazu müssen Eingaben, Szenenanalyse, Interaktionsplan, generierte Spezifikation pro Versuch, Validierungsfehler sowie die finale Ausgabe gespeichert werden.

\textbf{N3:} Die Implementierung des VSG soll so gestaltet werden, dass weitere Pipeline-Module und Validierungsregeln hinzugefügt werden können, ohne die Gesamtarchitektur grundlegend zu verändern.

\subsection*{Domänenspezifische Anforderungen}

\textbf{D1:} Der VSG erzeugt Vivian-Spezifikationsdateien, welche dem von Vivian erwarteten Dateiaufbau, den erlaubten Elementtypen, den Pflichtfeldern sowie zulässigen Wertebereichen entsprechen müssen.

\textbf{D2:} Der VSG darf vorhandene Geometrien und Szenenhierarchien nicht neu generieren oder ersetzen, sondern ausschließlich durch semantische und verhaltensbezogene Spezifikationen ergänzen.

\textbf{D3:} Die erzeugten Bezeichner für Interaktionselemente, Visualisierungselemente, Zustände, Übergänge und Objekt-Referenzen müssen eindeutig sein.

\textbf{D4:}
Jede Referenz in Zuständen, Übergängen, Events, Guards und Elementdefinitionen von Interaktionselementen sowie Visualisierungselementen muss auf ein vorhandenes Szenenobjekt oder ein vorhandenes Element der Spezifikation verweisen.

\textbf{D5:} Die generierte Spezifikation gilt als erfolgreich erzeugt, wenn sie in Unity mit dem Vivian Framework importiert und ohne Validierungsfehler ausgeführt werden kann.

% \begin{itemize}
% \item \textbf{Z1:} Entwicklung einer mehrstufigen Agenten-Architektur zur Zerlegung der Spezifikationsgenerierung in klar abgegrenzte Teilaufgaben.
% \item \textbf{Z2:} Entwicklung eines Systems zur automatisierten Erzeugung Vivian-kompatibler Spezifikationen aus Szenendaten und Sprache
% \item \textbf{Z3:} Integration von Validierungs- und Selbstkorrekturmechanismen zur Verbesserung syntaktischer und semantischer Konsistenz.
% \item \textbf{Z4:} Einbindung eines Human-in-the-Loop-Schritts zur Prüfung und Korrektur unklarer Szeneninterpretationen.
% \item \textbf{Z5:} Nutzer*innen sollen gewünschte Interaktionen durch Eingabe natürlicher Sprache spezifizieren können.
% \end{itemize}

\begin{table}[h]
    \centering
    \caption{Übersicht der Anforderungen an den VSG}
    \label{tab:anforderungen}
    \begin{tabular}{ll}
        \toprule
        \textbf{\#} & \textbf{Anforderung} \\
        \midrule
        A1          & XXX                  \\
        A2          & XXX                  \\
        A3          & XXX                  \\
        A4          & XXX                  \\
        A5          & XXX                  \\
        \midrule
        N1          & XXX                  \\
        N2          & XXX                  \\
        N3          & XXX                  \\
        \midrule
        D1          & XXX                  \\
        D2          & XXX                  \\
        D3          & XXX                  \\
        D4          & XXX                  \\
        D5          & XXX                  \\
        \bottomrule
    \end{tabular}
\end{table}

\section{Spezifikationsgenerierung für das Vivian Framework}
\section{Konzeption der Pipeline}
\subsection{Scene Analyzer}
\subsection{Interaction Planner}
\subsection{Specification Generator}
\subsection{Consistency Reviewer}
\subsection{Validator}